\clearpage
\section{Actual filling through expansion and unary term contexts}
\label{sec:copying-expansion-plug}
\termdef{expansion-replacement-interchange}
This section proves a structural part of the missing graph-filling bridge.
Replacing an edge within a component of a finite expansion agrees, up to
explicit graph isomorphism, with replacing that edge in the component before
expansion. The result applies to the original serial constructor and hence
to the full unary-operation clause of Figure~16, including its distributor
suffix. For a unary term context we obtain an actual addressed-filling law
whose target is the unchanged original term translation.

The result holds for arbitrary finite variable and atom environments and
arbitrary argument terms, including terms containing boxes. It proves a
unary context with its hole at the outer level; it does not yet prove
filling below an enclosing abstraction or inside a parallel argument pair.
Those cases and their common-allowance padding remain in Appendix~B.1,
item~3. No execution hypothesis is needed for the structural isomorphisms
proved here.

\subsection{The replacement--expansion isomorphism}
Let $C$ be a finite pattern and $F_b$ its assigned component at edge $b$.
Choose an edge $a$ of $C$, an edge $e$ of $F_a$, and a graph $H$ with the
interface of $e$. Write $F'_a=F_a[e:=H]$ and $F'_b=F_b$ for $b\ne a$.
The theorem constructs
\[
  (C[F])[(a,e):=H]\;\cong\; C[F'].
\]
It assumes no non-self-loop condition and no hole-free or binding-free
condition: it is a theorem about finite typed graph construction itself.


The changed family is defined by testing whether its component index is
$a$. At $a$ it is ordinary shallow replacement; elsewhere it is the original
component. Equality transport preserves the dependent interface type when
the index is identified with $a$. The two evaluation lemmas reduce the
selected and non-selected cases, respectively.

\lean{Graph/ExpansionPlug}{14}{18}{ExpansionPlug.family}


\lean{Graph/ExpansionPlug}{19}{21}{ExpansionPlug.family_same}


\lean{Graph/ExpansionPlug}{22}{25}{ExpansionPlug.family_other}


The source graph's internal carrier has the form
$(V_C+\sum_b V_{F_b})+V_H$. The target carrier is
$V_C+\sum_b V_{F'_b}$. Map each old pattern vertex to itself. A vertex in a
component different from $a$ keeps its tag. A vertex in $F_a$ becomes an
old vertex of $F_a[e:=H]$, and a vertex of $H$ becomes a newly inserted vertex
of that same component. The inverse splits these cases in reverse. The two
inverse laws follow by case analysis, using equality transport to identify
the selected component's internal carrier with $V_{F_a}+V_H$.

\lean{Graph/ExpansionPlug}{26}{36}{ExpansionPlug.vertexTo}


\lean{Graph/ExpansionPlug}{37}{47}{ExpansionPlug.vertexFrom}


\lean{Graph/ExpansionPlug}{48}{73}{ExpansionPlug.vertices}


The source edge carrier consists of all tagged component edges except
$(a,e)$, together with $H$'s edges. For a surviving edge in component $a$,
its inequality with $(a,e)$ proves that its local edge differs from $e$;
place it in the retained part of $F_a[e:=H]$. Other component edges are
unchanged, and new edges receive component tag $a$ and the inserted tag.
The inverse removes those tags. The inverse laws preserve both the dependent
component index and the retained-edge inequality; proof irrelevance removes
any difference between witnesses of the same inequality.

\lean{Graph/ExpansionPlug}{74}{84}{ExpansionPlug.edgeTo}


\lean{Graph/ExpansionPlug}{85}{97}{ExpansionPlug.edgeFrom}


\lean{Graph/ExpansionPlug}{98}{133}{ExpansionPlug.edges}


Transporting a graph vertex across an equality of graphs fixes all external
ports and transports only its internal carrier. Equality elimination proves
this fact. The vertex regrouping fixes every original pattern vertex,
including pattern boundary vertices. An unchanged component's vertex has
exactly the same attachment to the pattern before and after regrouping.
An old vertex of the selected component becomes a retained vertex in its
replacement graph, again with the same attachment.

For a newly inserted vertex, the internal case is the defining carrier map.
An input port attaches to the selected edge's corresponding source vertex;
an output attaches to its target vertex. The selected-component lemma,
together with expansion's ordered source and target list definitions,
shows that these attachments agree on both sides. Thus the maps preserve
actual incidence, not just counts or labels.

\lean{Graph/ExpansionPlug}{134}{139}{ExpansionPlug.cast_vertex}


\lean{Graph/ExpansionPlug}{140}{145}{ExpansionPlug.map_host}


\lean{Graph/ExpansionPlug}{146}{157}{ExpansionPlug.map_other}


\lean{Graph/ExpansionPlug}{158}{170}{ExpansionPlug.map_component}


\lean{Graph/ExpansionPlug}{171}{185}{ExpansionPlug.map_inserted}


Equality transport preserves internal labels and edge labels. It transports
source and target lists by mapping the corresponding vertex transport over
the original list. All four statements follow by eliminating the graph
equality; the incidence statements then reduce to mapping the identity
over a list. These lemmas expose the dependent casts required by the
changed family without changing graph structure.

\lean{Graph/ExpansionPlug}{186}{188}{ExpansionPlug.cast_internalLabel}


\lean{Graph/ExpansionPlug}{189}{191}{ExpansionPlug.cast_edgeLabel}


\lean{Graph/ExpansionPlug}{192}{195}{ExpansionPlug.cast_source}


\lean{Graph/ExpansionPlug}{196}{200}{ExpansionPlug.cast_target}


Three equations describe the edge map: one for retained edges of the selected
component, one for edges of other components, and one for newly inserted edges. They reduce directly
from the edge-map definition and the component-index equality test.

The final graph isomorphism uses the proved vertex and edge bijections.
Internal labels agree in each carrier case, and the three edge equations
reduce edge-label preservation to the transport lemmas. For source incidence,
a retained selected-component edge uses the selected-component vertex lemma;
an unchanged edge uses the other-component lemma; a new edge uses the
inserted-vertex lemma. Mapping these pointwise equalities over the entire
source list proves preservation of its order as well as its entries.
The target proof is identical with target lists. This completes the
replacement--expansion isomorphism with every interface and incidence law
checked by Lean.

\lean{Graph/ExpansionPlug}{201}{207}{ExpansionPlug.edgeTo_same}


\lean{Graph/ExpansionPlug}{208}{213}{ExpansionPlug.edgeTo_other}


\lean{Graph/ExpansionPlug}{214}{220}{ExpansionPlug.edgeTo_inserted}


\lean{Graph/ExpansionPlug}{221}{291}{ExpansionPlug.iso}


\subsection{Component isomorphisms and serial assembly}

A second generic construction assembles component-wise isomorphisms under
a common pattern. It permits a supplied relation on edge labels, so it also
applies to recursive graph isomorphism, where corresponding box bodies may
be isomorphic rather than literally equal.

Keep each pattern vertex fixed and apply the relevant component's vertex
bijection inside each tagged component. Apply its edge bijection inside
that component as well. The inverse laws are the component inverse laws.
Pattern boundary attachment is fixed, and component insertion commutes with
these maps by the internal/input/output cases. The resulting label-related
isomorphism inherits the internal-label and edge-relation proofs pointwise.
Its ordered incidence equations follow by mapping the insertion equality
over the component incidence lists and using the component isomorphism's
source or target equation.

\lean{Graph/ExpansionIso}{13}{23}{ExpansionIso.vertices}


\lean{Graph/ExpansionIso}{24}{29}{ExpansionIso.edges}


\lean{Graph/ExpansionIso}{30}{33}{ExpansionIso.map_host}


\lean{Graph/ExpansionIso}{34}{43}{ExpansionIso.map_insert}


\lean{Graph/ExpansionIso}{44}{64}{ExpansionIso.iso}


Serial composition is the original expansion of a two-slot pattern.
A left-component occurrence is tagged by the left Boolean slot, and a
right-component occurrence by the right slot. The edge maps are defined
at each allowance, respecting the copying alphabet's base and successor
representations. Both mapped occurrences retain exactly their original
edge labels, and hence their ordered interfaces.

\lean{Copying/SerialPlug}{13}{18}{SerialPlug.leftEdge}


\lean{Copying/SerialPlug}{19}{23}{SerialPlug.rightEdge}


\lean{Copying/SerialPlug}{24}{26}{SerialPlug.leftEdge_label}


\lean{Copying/SerialPlug}{27}{30}{SerialPlug.rightEdge_label}


To replace an edge on the left, apply the generic replacement--expansion
isomorphism with the selected slot equal to the left slot. The changed
family is pointwise equal to the serial filling family formed from the
replaced left graph and the untouched right graph. The two Boolean cases
prove that equality. Interface transport aligns the inserted graph with
the occurrence's retained label. The right-component theorem uses the
same argument with the other slot selected.

Finally, the generic component-isomorphism construction applied to these
two slots proves recursive isomorphism congruence for serial assembly.
These are laws for the original expanded graph, with no replacement serial
representation or additional semantic hypothesis.

\lean{Copying/SerialPlug}{31}{50}{SerialPlug.plug_left}


\lean{Copying/SerialPlug}{51}{69}{SerialPlug.plug_right}


\lean{Copying/SerialPlug}{70}{75}{SerialPlug.congr}


\subsection{The actual unary term-context filling law}
\termdef{copying-unary-filling-law}


The unary translation first serially connects its operation corolla to the
argument graph, then appends the prescribed single-use environment
distributor. Follow an argument occurrence through the right slot of the
first composition and then the left slot of the second. Its label remains
the argument edge's label by the two serial label lemmas.

For replacement at that occurrence, use the left serial replacement theorem
to move replacement past the distributor suffix. Use the right serial
replacement theorem to move it past the operation prefix. Lift the latter
isomorphism through the unchanged suffix by serial congruence and compose
the two isomorphisms. This gives the literal unary replacement law for any
argument edge and any fitting replacement graph. The corresponding unary
congruence theorem follows by applying serial congruence twice, fixing the
operation and distributor graphs.

\lean{Copying/UnaryPlug}{11}{15}{UnaryPlug.unaryEdge}


\lean{Copying/UnaryPlug}{16}{20}{UnaryPlug.unaryEdge_label}


\lean{Copying/UnaryPlug}{21}{45}{UnaryPlug.unary_plug}


\lean{Copying/UnaryPlug}{46}{51}{UnaryPlug.unary_congr}


Specialise the argument graph to the annotated hole corolla. Following its
single edge through the unary clause gives an explicit shallow site at the
hole's term-and-environment interface.

To certify that particular site, first prove a general occurrence fact.
If a graph has exactly one specified hole and an outer edge is labelled by
that hole, the outer edge is the selected occurrence: otherwise the
inventory's other-edge condition would declare the hole label hole-free,
a contradiction. Its shallow site is therefore sole. Applying this fact to
the existing unary single-hole inventory certifies the explicit address
used in the filling theorem. Thus the theorem does not rely on a different
existentially chosen hole address.

\lean{Copying/UnaryPlug}{52}{58}{UnaryPlug.unaryHoleSite}


\lean{Copying/UnaryPlug}{59}{74}{UnaryPlug.sole_of_single_hole}


\lean{Copying/UnaryPlug}{75}{88}{UnaryPlug.unaryHoleSite_sole}


Fill this actual site with an arbitrary graph $G$ at the annotated
interface. The unary replacement law identifies the result with the unary
translation of the hole corolla filled by $G$. The generic corolla left-unit
isomorphism identifies that filled corolla with $G$ itself. Unary congruence
lifts this isomorphism through the unchanged operation and distributor
construction. Writing $\mathcal{T}$ for the original translation, composing proves
\[
  \mathcal{T}(o([-]))[\mathcal{T}(t)]
  \;\cong\;\mathcal{T}(o(t))
\]
when $G$ is the original translation of $t$, with the hole allowance equal
to $t$'s translation depth. The final declaration states precisely this
specialisation using actual addressed replacement and the original
\texttt{Translation.term} on the right. All unary operations are covered,
including reference creation and dereference, for arbitrary scopes and
argument terms. This is a structural translation theorem, not a theorem
about executing either store operation.

\lean{Copying/UnaryPlug}{89}{105}{UnaryPlug.unary_hole_filling}


\lean{Copying/UnaryPlug}{106}{109}{UnaryPlug.unary_term_filling}


The full build and all three axiom audits pass. The dependency chain now
contains a proved unary filling case:
replacement--expansion interchange, serial replacement and congruence,
unary replacement, then corolla elimination. The general context induction
still needs replacement below boxes, the parallel argument construction,
and its padding compatibility. Term-context congruence and the copying
execution's representation laws remain separate obligations. The displayed
programs in Proposition~12.2 and every existing translation clause are unchanged.
